


We introduce GELATO, a novel approach to constructing multimodal embedding models by connecting frozen pre-trained modality-specific encoders directly to a frozen text embedding model via compact and easily trained projectors.
The result of this research, the \omniBase{} model suite, is also presented. 
These models add vision and audio to the Jina Embeddings v5 Text models, yielding a competitive set of models for broad cross-modality applications. 
Using GELATO, text-only embedding models that were never trained on vision or audio can be extended to photos, documents, video, speech, music, and sounds by training a single projector layer per modality while preserving text-only performance.

\omniSmall{} is the best-performing open-weight embedding model below $2$B parameters that supports text, audio, images, and video.
Against a baseline of comparable models, including modality-specific and VLM-derived embedders, it is particularly strong on visual document retrieval.
\omniSmall{} and \omniNano{} extend completely different text embedding models with different backbone architectures, suggesting that GELATO is an extensible strategy with broad application, outside of the \omniBase{} suite and for additional modalities.
This is a potential subject for future research.

The ablations suggest that projector-only alignment can serve as a compatibility-preserving initialization for rich multimodal training.
Future work will investigate the choice of non-text encoders, which is inadequately explored in this paper.
Furthermore, an investigation of training options under different conditions is indicated, like jointly training projectors for multiple modalities together.
We also note that \omniBase{} is comparatively closer to the baselines on moment retrieval than on the other video sub-tasks, but overall video performance remains weak.
We hope to improve performance in this area in future models.
